% Sumário executivo — elemento pré-textual do modo corporativo
% (\tipodocumento{corporativo} no main.tex).
%
% Diferente do resumo acadêmico, o sumário executivo é voltado a quem decide:
% abra com o contexto e a pergunta de negócio, apresente os principais achados
% e termine com as recomendações/próximos passos. Evite jargão técnico e
% priorize números e conclusões acionáveis. Sugestão de 1 a 2 páginas.

% [Substitua pelo conteúdo real do relatório]
Este relatório apresenta uma análise exploratória dos dados comerciais da
commodity em estudo ao longo do período analisado, com o objetivo de subsidiar
a tomada de decisão. A partir do tratamento e da consolidação dos dados,
identificaram-se os principais padrões de preço, volume e sazonalidade, bem
como os fatores que mais influenciaram o comportamento do mercado no intervalo
considerado.

Entre os achados, destacam-se: (i) a tendência geral observada no período;
(ii) os pontos de inflexão relevantes e suas prováveis causas; e (iii) os
segmentos ou regiões com maior potencial. Com base nessas evidências,
recomenda-se o conjunto de ações descrito ao final deste documento, cujo
detalhamento metodológico e analítico encontra-se nos capítulos seguintes.

Este texto é genérico e destina-se a demonstrar a formatação do template no
modo corporativo, devendo ser substituído pelo conteúdo real do relatório.

\palavraschave{Análise exploratória de dados. Commodity. Inteligência de mercado. Relatório corporativo.}
